problem:  
Let \((M^n,g)\) be a complete, oriented, asymptotically flat spin manifold of dimension \(n\ge 3\) with scalar curvature \(R_g\ge 0\) and ADM mass \(m_{\mathrm{ADM}}(g)\).  
Assume there exists a smooth, nontrivial spinor field \(\psi\) on \(M\) satisfying the **twistor equation**
\[
\nabla_X \psi + \frac{1}{n} X \cdot D\psi = 0 \quad \text{for all } X\in TM,
\]
and having the decay property \(|\psi| = O(r^{-\alpha})\) for some \(\alpha > (n-2)/2\) in asymptotic coordinates at spatial infinity.  

**Conjecture:**  
Under these assumptions, the ADM mass must vanish, \(m_{\mathrm{ADM}}(g)=0.\)  
By the rigidity part of the positive mass theorem, this would in particular imply that \((M,g)\) is isometric to \((\mathbb{R}^n,g_{\mathrm{eucl}})\).  

A stronger but related question is whether such existence of a decaying twistor spinor forces \((M,g)\) to be globally *conformally flat*, thus exhibiting a “conformal–spinorial rigidity” extending Witten’s spinorial framework beyond harmonic spinors.

Why is it a "good" problem:  
This conjecture proposes a precise and unexplored extension of the spinorial rigidity mechanism underlying the positive mass theorem. In Witten’s approach, a harmonic spinor satisfying \(D\psi = 0\) yields a boundary integral directly expressing the ADM mass, leading to its nonnegativity. Twistor spinors, however, satisfy a conformally covariant first-order system for which the Weitzenböck formula introduces additional curvature and divergence terms. The central analytic challenge—and the potential novelty—lies in whether these extra terms can still yield a nonnegative boundary quantity or enforce vanishing mass via decay constraints.  

**Profound insight:**  
It asks whether the ADM mass, a global gravitational invariant, is directly constrained by the existence of conformally covariant spinorial fields, suggesting a deeper conformal nature of mass positivity and rigidity.  

**Bridge between fields:**  
It unites spin geometry, conformal geometry, spectral theory of the Dirac operator, and general relativity. Progress would demand reconciling the Weitzenböck identities for twistor spinors with the asymptotic analysis of Witten’s boundary integral—an intersection previously unexplored.  

**Extreme simplicity:**  
Despite involving sophisticated structures, the conjecture rests on a single equation (the twistor equation) and one asymptotic condition, analogous in form to Witten’s harmonic spinor—making its statement concise and conceptually clean.  

**Nontrivial and open:**  
The classification of manifolds carrying twistor spinors is well understood in closed or Einstein settings but virtually unknown for asymptotically flat manifolds. If proven, the conjecture would either reveal that decaying twistor spinors cannot exist except in the Euclidean case or uncover a new analytic mechanism relating conformal spinorial structures to gravitational mass.